%!TEX program = xelatex
%!TEX encoding = ISO-8859-1
%!TEX program = xelatex
% !TEX encoding = ISO-8859-1
\def\ano{2019}
\def\semestre{2}
\def\disciplina{\'Algebra 1}
\def\turma{C}

\documentclass[12pt]{exam}

\usepackage{caption}
\usepackage{amssymb}
\usepackage{amsmath,amsfonts,amsthm,amstext}
\usepackage[brazil]{babel}
% \usepackage[latin1]{inputenc}
\usepackage{graphicx}
\graphicspath{{/home/jfreitas/GitHub_Repos/Algebra-1/Pictures/}{D:/OneDrive - unb.br/imagens-latex/}}
\usepackage{enumitem}
\usepackage{multicol}
\usepackage{answers}
\usepackage{tikz,ifthen}
\usetikzlibrary{lindenmayersystems}
\usetikzlibrary[shadings]
\Newassociation{solucao}{Solution}{ans}
\newtheorem{exercicio}{}

\setlength{\topmargin}{-1.0in}
\setlength{\oddsidemargin}{0in}
\setlength{\textheight}{10.1in}
\setlength{\textwidth}{6.5in}
\setlength{\baselineskip}{12mm}

\extraheadheight{0.7in}
\firstpageheadrule
\runningheadrule
\lhead{
        \begin{minipage}[c]{1.7cm}
        \includegraphics[width=1.7cm]{unb.pdf}
        \end{minipage}%
        \hspace{0pt}
        \begin{minipage}[c]{4in}
          {Universidade de Bras{\'\i}lia} --
          {Departamento de Matem{\'a}tica}
\end{minipage}
\vspace*{-0.8cm}
}
% \chead{Universidade de Bras{\'\i}lia - Departamento de Matem{\'a}tica}
% \rhead{}
% \vspace*{-2cm}

\extrafootheight{.5in}
\footrule
\lfoot{\disciplina\ - \semestre$^o$/\ano}
\cfoot{}
\rfoot{P\'agina \thepage\ de \numpages}

\newcounter{exercicios}
\renewcommand{\theexercicios}{\arabic{exercicios}}

\newenvironment{questao}[1]{
\refstepcounter{exercicios}
\ifx&#1&
\else
   \label{#1}
\fi
\noindent\textbf{Exerc{\'\i}cio {\theexercicios}:}
}

\newcommand{\resp}[1]{
\noindent{\bf Exerc{\'\i}cio #1: }}
\newcommand{\nsub}{\varsubsetneq}
\newcommand{\vaz}{\emptyset}
\newcommand{\im}{{\rm Im\,}}
\newcommand{\sub}{\subseteq}
\newcommand{\n}{\mathbb{N}}
\newcommand{\z}{\mathbb{Z}}
\newcommand{\rac}{\mathbb{Q}}
\newcommand{\real}{\mathbb{R}}
\newcommand{\complex}{\mathbb{C}}
\newcommand{\cp}[1]{\mathbb{#1}}
\newcommand{\ch}{\mbox{\textrm{car\,}}\nobreak}
\newcommand{\vesp}[1]{\vspace{ #1  cm}}
\newcommand{\compcent}[1]{\vcenter{\hbox{$#1\circ$}}}
\newcommand{\comp}{\mathbin{\mathchoice
{\compcent\scriptstyle}{\compcent\scriptstyle}
{\compcent\scriptscriptstyle}{\compcent\scriptscriptstyle}}}

\begin{document}

\begin{center}
{\Large\bf \disciplina\ - Turma \turma\ -- \semestre$^{o}$/\ano} \\ \vspace{9pt} {\large\bf
$3^{\underline{o}}$ Teste - M\'odulo 1 - Resolu\c{c}\~ao}\\
\vspace{9pt} Prof. Jos{\'e} Ant{\^o}nio O. Freitas
\end{center}
\hrule

\vspace{.6cm}

\questao Seja $m \in \z$, $m > 1$. Defina $R \sub \z \times \z$ por
\[
	R = \{ (x, y) \in \z \times \z \mid x - y = mk,\ k \in \z\}.
\]

Mostre que $R$ \'e uma rela\c{c}\~ao de equival\^encia sobre $\z$.

\noindent\textbf{Solu\c{c}\~ao:} Para mostrar que $R$ \'e rela\c{c}\~ao de equival\^encia sobre um conjunto $A \ne \vaz$ precisamos mostrar que:
\begin{enumerate}[label={\roman*})]
	\item Para todo $x \in A$, $(x, x) \in R$.
	\item Se $(x, y) \in R$, ent\~ao $(y, x) \in R$.
	\item Se $(x, y) \in R$ e $(y, z) \in R$, ent\~ao $(x, z) \in R$.
\end{enumerate}

De fato,
\begin{enumerate}[label={\roman*})]
	\item Para $x \in \z$ temos $x - x = 0 = m\cdot 0$, logo $(x, x) \in R$.
	\item Se $(x, y)\in R$ ent\~ao, para algum $k \in \z$ temos $x - y  = mk$. Assim
	\begin{align*}
		y - x = -(x - y) = -(mk) = m(-k)
	\end{align*}
	logo $(y, x) \in R$.
	\item Se $(x, y) \in R$ e $(y, z) \in R$ ent\~ao, existem $k$, $l \in \z$ tais que
	\begin{align*}
		x - y &= mk\\
		y - z &= ml.
	\end{align*}
	Assim
	\begin{align*}
		x - z = (x - y) + (y - z) = km + lm = m(k + l)
	\end{align*}
	e ent\~ao $(x, z) \in R$.
\end{enumerate}

Portanto, $R$ \'e uma rela\c{c}\~ao de equival\^encia sobre $\z$ como quer{\'\i}amos.

\vspace{.5cm}

\questao Seja $\sim$ a rela\c{c}\~ao sobre $\rac$ definida da seguinte forma
\[
	x\sim y \mbox{ quando } x - y \in \z.
\]

Mostre que $R$ \'e uma rela{\c c}{\~a}o de equival{\^e}ncia sobre $\rac$. Encontre $\overline{1}$ e $\overline{1/2}$.


\noindent\textbf{Solu\c{c}\~ao:}

Primeiro dizer que $x - y \in \z$, significa que existe um $k \in \z$ tal que $x - y = k$. Assim:
\begin{enumerate}[label={\roman*})]
	\item Para todo $x \in \rac$, $x - x = 0 \in \rac$. Logo $x \sim x$.
	\item Se $x \sim y$, ent\~ao existe $k \in \z$ tal que $x - y = k$. Da{\'\i}
	\[
		y - x = -(x - y) = -k
	\]
	e como $-k \in \z$, ent\~ao $y \sim x$.
	\item Se $x \sim y$ e $y \sim z$, ent\~ao existem $k$, $l \in \z$ tais que
	\begin{align*}
		x - y &= k\\
		y - z &= l.
	\end{align*}
	Logo
	\[
		x - z = (x - y) + (y - z) = k + l
	\]
	e como $k + l \in \z$, ent\~ao $x \sim z$.
\end{enumerate}

Portanto, $\sim$ \'e uma rela\c{c}\~ao de equival\^encia sobre $\rac$ como quer{\'\i}amos.

Agora,
\begin{align*}
	\overline{1} &= \{y \in \rac \mid 1 \sim y\} \\&= \{y \in \rac \mid 1 - y \in \z\} \\&= \{y \in \rac \mid 1 - y = k,\ k \in \z\}\\&= \{y \in \rac \mid y = 1 - k, k \in \z\} \\&= \{k + 1 \mid k \in \z\}
\end{align*}
\begin{align*}
	\overline{1/2} &= \{y \in \rac \mid 1/2 \sim y\} \\&= \{y \in \rac \mid 1/2 - y \in \z\} \\&= \{y \in \rac \mid 1/2 - y = k,\ k \in \z\} \\&= \{y \in \rac \mid y = 1/2 - k, k \in \z\} \\&= \left\{\dfrac{1}{2} - k \mid k \in \z\right\}
\end{align*}
\end{document}